%% ========================================================================
%% Chapter 12: Manajemen Layanan (Services)
%% ========================================================================

\chapter{Manajemen Layanan (Services)}
\label{ch:service-management}

Bab ini membahas bagaimana Linux mengelola layanan sistem dan layanan aplikasi,
mulai dari konsep \textit{init system}, penggunaan \cmd{systemctl}, pembuatan
\textit{unit file}, pengaturan dependensi, pemantauan log, hingga troubleshooting
layanan yang gagal start. Fokus utama bab ini adalah \textbf{systemd} karena ia
menjadi \textit{init system} standar pada Ubuntu Server 22.04 dan mayoritas
distribusi Linux modern.

\textbf{Tujuan Praktikum}

Setelah menyelesaikan bab ini, pembaca diharapkan mampu:

\begin{enumerate}
    \item menjelaskan peran sistem init di Linux dan evolusinya;
    \item mengelola layanan dengan \cmd{systemctl};
    \item membaca, membuat, dan menyesuaikan \textit{unit file} systemd;
    \item memahami dependensi dan urutan start layanan;
    \item memantau status layanan dan log dengan \cmd{journalctl};
    \item mengonfigurasi layanan jaringan umum seperti SSH, Apache, dan MySQL;
    \item mengotomatisasi start-up dan shutdown layanan; dan
    \item melakukan troubleshooting saat layanan gagal berjalan.
\end{enumerate}

\begin{notebox}
Seluruh contoh di bab ini diasumsikan berjalan di Ubuntu Server 22.04. Sebagian
besar perintah yang memodifikasi layanan membutuhkan \cmd{sudo}. Gunakan direktori
kerja terpisah untuk eksperimen file unit kustom:
\begin{lstlisting}[language=bash, caption={Menyiapkan direktori kerja praktikum layanan}]
mkdir -p ~/praktikum-os/week12-services
cd ~/praktikum-os/week12-services
\end{lstlisting}
\end{notebox}

%% ========================================================================
\section{Pengenalan Sistem Init}
\label{sec:init-system}
%% ========================================================================

Sistem init adalah proses pertama di ruang pengguna yang dijalankan kernel setelah
boot selesai memuat sistem dasar. Proses ini biasanya memiliki PID 1, dan bertugas
melanjutkan proses boot, memulai layanan latar belakang, menjaga layanan tetap
berjalan, serta menangani shutdown atau reboot sistem.

Tanpa sistem init, kernel memang dapat berjalan, tetapi sistem tidak akan otomatis
menjalankan layanan penting seperti SSH, logging, network manager, atau database.
Karena itu, init system adalah penghubung utama antara kernel dan layanan pengguna.

\subsection{Evolusi Init di Linux}

Secara historis Linux mengenal beberapa pendekatan init system:

\begin{description}
    \item[SysV init] Model klasik berbasis skrip shell di direktori seperti
        \dir{/etc/init.d/}. Ia sederhana dan mudah dipahami, tetapi urutan start
        bersifat relatif linear dan manajemen dependensinya terbatas.
    \item[Upstart] Pendekatan berbasis event yang pernah digunakan Ubuntu sebelum
        migrasi ke systemd. Upstart memperbaiki sebagian keterbatasan SysV init,
        tetapi adopsinya tidak menjadi standar lintas distribusi.
    \item[systemd] Init system modern yang berbasis \textit{unit}. Ia menyediakan
        start paralel, manajemen dependensi yang eksplisit, socket activation,
        logging terintegrasi, dan alat administrasi yang seragam.
\end{description}

\begin{table}[htbp]
    \centering
    \caption{Perbandingan pendekatan init system di Linux}
    \label{tab:init-systems}
    \begin{tblr}{
        colspec = {lX[l]X[l]X[l]},
        width = \linewidth,
        hlines,
        vlines,
        row{1} = {font=\bfseries},
        cells = {font=\small}
    }
    \textbf{Aspek} & \textbf{SysV init} & \textbf{Upstart} & \textbf{systemd} \\
    Model utama & Skrip shell berurutan & Event-based jobs & Unit-based service manager \\
    Dependensi & Terbatas, implisit & Lebih baik dari SysV & Eksplisit dan kaya \\
    Kecepatan boot & Cenderung serial & Lebih fleksibel & Start paralel, lebih cepat \\
    Logging & Terpisah & Terpisah & Terintegrasi dengan journal \\
    Status saat ini & Legacy / kompatibilitas & Hampir tidak dipakai & Standar modern \\
    \end{tblr}
\end{table}

Di Ubuntu Server 22.04, perintah utama administrasi layanan terpusat pada
\cmd{systemctl} untuk kontrol layanan dan \cmd{journalctl} untuk log.

\begin{lstlisting}[language=bash, caption={Memeriksa init system yang digunakan}]
ps -p 1 -o pid,comm,args=
systemctl --version
\end{lstlisting}

Jika PID 1 menunjukkan \texttt{systemd}, berarti seluruh manajemen layanan modern
pada sistem tersebut dikendalikan oleh systemd.

%% ========================================================================
\section{Mengelola Layanan dengan systemctl}
\label{sec:systemctl-services}
%% ========================================================================

\cmd{systemctl} adalah antarmuka utama untuk mengelola unit systemd. Walau systemd
tidak hanya mengelola \textit{service}, dalam praktik administrasi sehari-hari
unit yang paling sering disentuh memang unit layanan.

\subsection{Perintah Dasar systemctl}

\begin{lstlisting}[language=bash, caption={Perintah dasar mengelola layanan}]
sudo systemctl start ssh
sudo systemctl stop ssh
sudo systemctl restart ssh
sudo systemctl reload ssh
sudo systemctl status ssh
sudo systemctl is-active ssh
sudo systemctl is-enabled ssh
\end{lstlisting}

Makna perintah-perintah di atas:

\begin{itemize}
    \item \cmd{start}: menjalankan layanan jika sedang mati.
    \item \cmd{stop}: menghentikan layanan.
    \item \cmd{restart}: menghentikan lalu menjalankan ulang layanan.
    \item \cmd{reload}: memuat ulang konfigurasi tanpa memutus proses utama, jika
        layanan mendukungnya.
    \item \cmd{status}: menampilkan ringkasan status, PID utama, dan log terakhir.
    \item \cmd{is-active}: mengembalikan status singkat seperti \texttt{active}
        atau \texttt{inactive}.
    \item \cmd{is-enabled}: memeriksa apakah layanan otomatis start saat boot.
\end{itemize}

\begin{warningbox}
Jangan langsung mengasumsikan \cmd{restart} aman pada semua layanan. Untuk database,
web server produksi, atau aplikasi stateful, restart dapat memutus koneksi aktif.
Jika cukup, gunakan \cmd{reload} agar konfigurasi baru dimuat tanpa mematikan proses.
\end{warningbox}

\subsection{Enable, Disable, Mask, dan Unmask}

Selain status runtime, systemd juga mengelola apakah layanan akan berjalan otomatis
saat boot.

\begin{lstlisting}[language=bash, caption={Mengatur layanan saat boot}]
sudo systemctl enable ssh
sudo systemctl disable ssh
sudo systemctl enable --now apache2
sudo systemctl disable --now apache2
sudo systemctl mask cups
sudo systemctl unmask cups
\end{lstlisting}

Perbedaan istilah penting:

\begin{description}
    \item[enable] Membuat symlink agar unit dipanggil oleh target boot yang relevan.
    \item[disable] Menghapus symlink auto-start, tetapi layanan masih bisa dijalankan manual.
    \item[mask] Memblokir layanan sepenuhnya dengan mengarahkannya ke \file{/dev/null},
        sehingga bahkan dependensi lain tidak bisa men-start-nya tanpa \cmd{unmask}.
\end{description}

\begin{table}[htbp]
    \centering
    \caption{Perbedaan status pengelolaan layanan di systemd}
    \label{tab:service-states}
    \begin{tblr}{
        colspec = {lQ[5cm,l]X[l]},
        width = \linewidth,
        hlines,
        vlines,
        row{1} = {font=\bfseries},
        cells = {font=\small}
    }
    \textbf{Status} & \textbf{Makna} & \textbf{Dapat start manual?} \\
    active & Layanan sedang berjalan & Ya \\
    inactive & Layanan tidak berjalan & Ya \\
    enabled & Akan start otomatis saat boot & Ya \\
    disabled & Tidak start otomatis saat boot & Ya \\
    masked & Diblokir total oleh systemd & Tidak \\
    failed & Start gagal atau crash & Bergantung penyebab \\
    \end{tblr}
\end{table}

\subsection{Melihat Daftar Unit}

\begin{lstlisting}[language=bash, caption={Melihat unit systemd}]
systemctl list-units --type=service
systemctl list-units --type=service --state=running
systemctl list-unit-files --type=service
systemctl list-dependencies multi-user.target
\end{lstlisting}

Perintah-perintah ini berguna untuk audit layanan yang aktif, layanan yang tersedia,
dan hubungan dependensi pada target boot tertentu.

%% ========================================================================
\section{Membuat dan Menyesuaikan File Unit Sistem}
\label{sec:unit-files}
%% ========================================================================

Systemd mengelola berbagai jenis unit seperti \texttt{service}, \texttt{socket},
\texttt{timer}, \texttt{mount}, dan \texttt{target}. Untuk layanan aplikasi,
unit yang paling umum adalah \texttt{.service}.

\subsection{Struktur Unit File}

File unit biasanya berada di salah satu lokasi berikut:

\begin{itemize}
    \item \file{/lib/systemd/system/} atau \file{/usr/lib/systemd/system/}: unit bawaan paket.
    \item \file{/etc/systemd/system/}: override dan unit kustom lokal.
    \item \file{~/.config/systemd/user/}: unit level user.
\end{itemize}

Sebuah unit \texttt{service} sederhana memiliki tiga blok utama:

\begin{lstlisting}[language=bash, caption={Contoh unit file systemd sederhana}]
[Unit]
Description=Simple HTTP Demo Service
After=network.target

[Service]
Type=simple
ExecStart=/usr/bin/python3 -m http.server 9000
WorkingDirectory=/home/user/demo
Restart=on-failure

[Install]
WantedBy=multi-user.target
\end{lstlisting}

Penjelasan blok-bloknya:

\begin{description}
    \item[[Unit]] Metadata umum dan dependensi dasar.
    \item[[Service]] Cara proses dijalankan, dihentikan, dan dipantau.
    \item[[Install]] Menentukan bagaimana unit dihubungkan ke target saat di-enable.
\end{description}

\subsection{Field Penting pada [Service]}

\begin{table}[htbp]
    \centering
    \caption{Field penting dalam unit \texttt{service}}
    \label{tab:service-unit-fields}
    \begin{tblr}{
        colspec = {lX[l]X[l]},
        width = \linewidth,
        hlines,
        vlines,
        row{1} = {font=\bfseries},
        cells = {font=\small}
    }
    \textbf{Field} & \textbf{Contoh} & \textbf{Fungsi} \\
    Type & \texttt{simple} & Menentukan model proses utama layanan \\
    ExecStart & \texttt{/usr/bin/app} & Perintah utama untuk start layanan \\
    ExecStop & \texttt{pkill app} & Perintah saat layanan dihentikan \\
    WorkingDirectory & \texttt{/opt/app} & Direktori kerja proses \\
    User & \texttt{User=www-data} & Menjalankan layanan sebagai user tertentu \\
    Group & \texttt{Group=www-data} & Menjalankan layanan sebagai group tertentu \\
    Restart & \texttt{Restart=on-failure} & Kebijakan restart otomatis \\
    Environment & \texttt{Environment=PORT=8080} & Variabel lingkungan untuk proses \\
    \end{tblr}
\end{table}

\subsection{Membuat Unit Kustom}

Contoh berikut membuat layanan sederhana yang menjalankan server HTTP bawaan Python.

\begin{lstlisting}[language=bash, caption={Membuat unit file kustom}]
mkdir -p ~/praktikum-os/week12-services/demo-site
echo "Halo dari layanan systemd" > ~/praktikum-os/week12-services/demo-site/index.html

sudo tee /etc/systemd/system/demo-http.service > /dev/null <<'EOF'
[Unit]
Description=Demo HTTP Service
After=network.target

[Service]
Type=simple
User=$USER
WorkingDirectory=/home/$USER/praktikum-os/week12-services/demo-site
ExecStart=/usr/bin/python3 -m http.server 9000
Restart=on-failure

[Install]
WantedBy=multi-user.target
EOF
\end{lstlisting}

Setelah membuat atau mengubah unit file, systemd harus diminta membaca ulang definisinya:

\begin{lstlisting}[language=bash, caption={Memuat ulang unit systemd dan menjalankan layanan}]
sudo systemctl daemon-reload
sudo systemctl start demo-http
sudo systemctl status demo-http
curl http://localhost:9000
\end{lstlisting}

\begin{notebox}
Perintah \cmd{daemon-reload} tidak me-restart semua layanan. Ia hanya meminta systemd
membaca ulang definisi unit file dari disk.
\end{notebox}

\subsection{Override Tanpa Mengubah File Paket}

Untuk layanan bawaan paket, hindari mengedit unit bawaan paket secara langsung.
Gunakan mekanisme override:

\begin{lstlisting}[language=bash, caption={Membuat override unit file}]
sudo systemctl edit ssh
\end{lstlisting}

Perintah ini membuat file override di bawah
\file{/etc/systemd/system/} lalu menyimpan override sebagai
\file{ssh.service.d/override.conf}.
Dengan pendekatan ini, perubahan tidak hilang saat paket di-upgrade.

%% ========================================================================
\section{Manajemen Dependensi Layanan}
\label{sec:service-dependencies}
%% ========================================================================

Layanan di Linux jarang berdiri sendiri. Web application mungkin membutuhkan jaringan,
database, dan filesystem tertentu sebelum bisa start. Systemd menangani relasi ini
melalui dependensi dan aturan urutan eksekusi.

\subsection{After, Before, Wants, Requires}

Empat direktif yang paling penting adalah:

\begin{description}
    \item[After=] Menyatakan urutan start. Unit ini start setelah unit lain.
    \item[Before=] Kebalikan dari \texttt{After=}.
    \item[Wants=] Relasi lemah. Jika unit ini start, systemd juga mencoba start unit lain.
    \item[Requires=] Relasi kuat. Jika dependensi gagal, unit ini ikut gagal.
\end{description}

Contoh:

\begin{lstlisting}[language=bash, caption={Contoh dependensi pada unit file}]
[Unit]
Description=Aplikasi Web Internal
After=network.target mysql.service
Wants=network.target
Requires=mysql.service
\end{lstlisting}

Maknanya:

\begin{itemize}
    \item layanan baru harus menunggu jaringan dan MySQL start lebih dahulu;
    \item systemd akan mencoba memastikan network target aktif;
    \item jika MySQL gagal start, aplikasi ini juga dianggap gagal dijalankan.
\end{itemize}

\subsection{Target dan Runlevel Modern}

Systemd menggantikan konsep runlevel tradisional dengan \textit{target}.

\begin{table}[htbp]
    \centering
    \caption{Padanan runlevel lama dengan target systemd}
    \label{tab:runlevel-targets}
    \begin{tblr}{
        colspec = {lX[l]X[l]},
        width = \linewidth,
        hlines,
        vlines,
        row{1} = {font=\bfseries},
        cells = {font=\small}
    }
    \textbf{Runlevel lama} & \textbf{Target systemd} & \textbf{Makna} \\
    0 & \texttt{poweroff.target} & Shutdown sistem \\
    1 & \texttt{rescue.target} & Mode pemulihan / single-user \\
    3 & \texttt{multi-user.target} & Sistem multi-user non-grafis \\
    5 & \texttt{graphical.target} & Sistem grafis penuh \\
    6 & \texttt{reboot.target} & Reboot sistem \\
    \end{tblr}
\end{table}

Periksa target default sistem:

\begin{lstlisting}[language=bash, caption={Melihat dan mengubah target default}]
systemctl get-default
sudo systemctl set-default multi-user.target
sudo systemctl set-default graphical.target
\end{lstlisting}

%% ========================================================================
\section{Logging dan Monitoring}
\label{sec:service-logging}
%% ========================================================================

Systemd terintegrasi erat dengan journald. Alih-alih hanya bergantung pada file log
tradisional di \dir{/var/log/}, administrator dapat melihat log layanan langsung
melalui \cmd{journalctl}.

\subsection{Memantau Status Layanan}

\begin{lstlisting}[language=bash, caption={Memantau status layanan}]
systemctl status ssh
systemctl status apache2
systemctl --failed
systemctl show ssh --property=MainPID,SubState,ActiveEnterTimestamp
\end{lstlisting}

\cmd{systemctl --failed} sangat berguna saat audit cepat karena langsung menampilkan
semua unit yang gagal.

\subsection{Membaca Log dengan journalctl}

\begin{lstlisting}[language=bash, caption={Membaca log layanan dengan journalctl}]
journalctl -u ssh
journalctl -u ssh -n 50
journalctl -u ssh -f
journalctl -u apache2 --since "today"
journalctl -p err -b
\end{lstlisting}

Perintah-perintah penting:

\begin{itemize}
    \item \cmd{-u nama-unit}: hanya log untuk unit tertentu.
    \item \cmd{-n 50}: tampilkan 50 entri terakhir.
    \item \cmd{-f}: mode follow, seperti \cmd{tail -f}.
    \item \cmd{--since}: filter berdasarkan waktu.
    \item \cmd{-p err}: hanya log dengan prioritas error.
    \item \cmd{-b}: log sejak boot saat ini.
\end{itemize}

\begin{tipbox}
Saat layanan gagal start, kombinasi tercepat biasanya adalah: \cmd{systemctl status nama}
lalu \cmd{journalctl -u nama -n 50}. Dua perintah ini hampir selalu menjadi titik awal
troubleshooting.
\end{tipbox}

%% ========================================================================
\section{Konfigurasi Layanan Jaringan}
\label{sec:network-services}
%% ========================================================================

Bagian ini menggunakan tiga layanan umum sebagai contoh: SSH, Apache, dan MySQL.

\subsection{SSH}

SSH adalah layanan administrasi jarak jauh yang hampir selalu aktif di server Linux.

\begin{lstlisting}[language=bash, caption={Mengelola layanan SSH}]
sudo systemctl status ssh
sudo systemctl restart ssh
sudo systemctl enable ssh
sudo ss -tlnp | grep :22
\end{lstlisting}

Konfigurasi utamanya berada di \file{/etc/ssh/sshd\_config}. Setelah mengubah file ini,
gunakan \cmd{sshd -t} untuk memverifikasi sintaks sebelum restart:

\begin{lstlisting}[language=bash, caption={Validasi konfigurasi SSH}]
sudo sshd -t
sudo systemctl reload ssh
\end{lstlisting}

\subsection{Apache}

\begin{lstlisting}[language=bash, caption={Mengelola layanan Apache}]
sudo apt install apache2
sudo systemctl status apache2
sudo systemctl enable --now apache2
sudo apache2ctl configtest
sudo systemctl reload apache2
\end{lstlisting}

Port HTTP standar dapat diverifikasi dengan:

\begin{lstlisting}[language=bash, caption={Memeriksa port Apache}]
sudo ss -tlnp | grep :80
curl http://localhost
\end{lstlisting}

\subsection{MySQL}

\begin{lstlisting}[language=bash, caption={Mengelola layanan MySQL}]
sudo apt install mysql-server
sudo systemctl status mysql
sudo systemctl enable --now mysql
sudo mysqladmin ping
sudo ss -tlnp | grep :3306
\end{lstlisting}

Untuk database, verifikasi tidak cukup hanya melihat \cmd{active (running)}; selalu
cek juga apakah layanan benar-benar merespons koneksi aplikasi.

%% ========================================================================
\section{Automatisasi Layanan}
\label{sec:service-automation}
%% ========================================================================

Automatisasi layanan mencakup dua aspek: memastikan layanan start atau stop pada
waktu yang tepat, dan menjalankan tugas administratif berkala tanpa intervensi manual.

\subsection{Start-up dan Shutdown Otomatis}

Cara paling dasar adalah menggunakan \cmd{enable} dan \cmd{disable} seperti yang
sudah dibahas. Namun systemd juga mendukung hook start dan stop di unit file:

\begin{lstlisting}[language=bash, caption={Hook start/stop dalam unit file}]
[Service]
ExecStartPre=/usr/bin/test -f /etc/myapp/config.yml
ExecStart=/usr/local/bin/myapp
ExecStop=/usr/bin/pkill -f myapp
ExecStopPost=/usr/bin/logger "myapp berhenti"
\end{lstlisting}

\subsection{Timer sebagai Pengganti Cron untuk Tugas Tertentu}

Walau outline minggu ini berfokus pada layanan, penting mengenal \texttt{.timer}
karena ia adalah pasangan modern bagi layanan satu kali atau periodik.

\begin{lstlisting}[language=bash, caption={Contoh service dan timer sederhana}]
# /etc/systemd/system/cleanup-temp.service
[Unit]
Description=Bersihkan file sementara

[Service]
Type=oneshot
ExecStart=/usr/bin/find /tmp/myapp -type f -mtime +7 -delete

# /etc/systemd/system/cleanup-temp.timer
[Unit]
Description=Jalankan cleanup-temp setiap hari

[Timer]
OnCalendar=daily
Persistent=true

[Install]
WantedBy=timers.target
\end{lstlisting}

Aktivasi timer:

\begin{lstlisting}[language=bash, caption={Mengaktifkan timer systemd}]
sudo systemctl daemon-reload
sudo systemctl enable --now cleanup-temp.timer
systemctl list-timers
\end{lstlisting}

\begin{notebox}
\texttt{Persistent=true} membuat timer yang terlewat saat sistem mati dijalankan
setelah sistem hidup kembali, perilaku yang sering diinginkan untuk maintenance.
\end{notebox}

%% ========================================================================
\section{Troubleshooting}
\label{sec:service-troubleshooting}
%% ========================================================================

Troubleshooting layanan hampir selalu melibatkan empat langkah: periksa status,
baca log, verifikasi konfigurasi, lalu uji dependensi eksternal seperti port,
permission, atau file path.

\subsection{Pola Diagnostik Dasar}

\begin{lstlisting}[language=bash, caption={Pola troubleshooting layanan}]
systemctl status nama-layanan
journalctl -u nama-layanan -n 50
sudo ss -tlnp
ps aux | grep nama-layanan
\end{lstlisting}

Penyebab umum layanan gagal start:

\begin{itemize}
    \item file konfigurasi salah sintaks;
    \item binary atau path pada \texttt{ExecStart=} tidak ada;
    \item port yang dibutuhkan sudah dipakai proses lain;
    \item permission file atau direktori tidak memadai;
    \item dependensi seperti jaringan, database, atau mount point belum siap.
\end{itemize}

\subsection{Contoh Kasus Troubleshooting}

Misalnya sebuah layanan web gagal start karena port 8080 sudah digunakan:

\begin{lstlisting}[language=bash, caption={Mendeteksi konflik port}]
sudo systemctl start demo-http
sudo systemctl status demo-http
sudo journalctl -u demo-http -n 20
sudo ss -tlnp | grep :8080
\end{lstlisting}

Jika hasil \cmd{ss} menunjukkan proses lain memakai port yang sama, ada dua opsi:
ubah konfigurasi port layanan atau hentikan proses yang bentrok.

\begin{warningbox}
Jangan terburu-buru menjalankan \cmd{kill -9} pada proses yang tampak mengganggu.
Pastikan Anda memahami proses tersebut milik layanan apa, dan apakah penghentiannya
aman bagi sistem produksi.
\end{warningbox}

%%: Praktikum ---
\subsection*{Praktikum 10.1: Mengelola Layanan dengan systemctl}

\textbf{Tujuan:} memahami start, stop, restart, enable, disable, dan inspeksi log.

\textbf{Langkah-langkah:}
\begin{enumerate}
    \item Periksa status beberapa layanan umum:
\begin{lstlisting}[language=bash, caption={Memeriksa status layanan umum}]
systemctl status ssh
systemctl status cron
systemctl --failed
\end{lstlisting}

    \item Aktifkan satu layanan yang tersedia pada sistem Anda:
\begin{lstlisting}[language=bash, caption={Enable layanan saat boot}]
sudo systemctl enable ssh
sudo systemctl is-enabled ssh
\end{lstlisting}

    \item Lakukan restart dan baca log terbarunya:
\begin{lstlisting}[language=bash, caption={Restart dan inspeksi log layanan}]
sudo systemctl restart ssh
systemctl status ssh
journalctl -u ssh -n 20
\end{lstlisting}
\end{enumerate}

%%: Praktikum ---
\subsection*{Praktikum 10.2: Membuat Unit File Kustom}

\textbf{Tujuan:} membuat layanan sederhana dan menjalankannya di bawah systemd.

\textbf{Langkah-langkah:}
\begin{enumerate}
    \item Buat direktori dan halaman sederhana:
\begin{lstlisting}[language=bash, caption={Menyiapkan konten layanan kustom}]
mkdir -p ~/praktikum-os/week12-services/site
echo "<h1>Service Demo</h1>" > ~/praktikum-os/week12-services/site/index.html
\end{lstlisting}

    \item Buat unit file kustom:
\begin{lstlisting}[language=bash, caption={Membuat demo-http.service}]
sudo tee /etc/systemd/system/demo-http.service > /dev/null <<'EOF'
[Unit]
Description=Demo HTTP Service
After=network.target

[Service]
Type=simple
ExecStart=/usr/bin/python3 -m http.server 9000 --directory /home/$USER/praktikum-os/week12-services/site
Restart=on-failure

[Install]
WantedBy=multi-user.target
EOF
\end{lstlisting}

    \item Muat ulang daemon dan jalankan layanannya:
\begin{lstlisting}[language=bash, caption={Menjalankan layanan kustom}]
sudo systemctl daemon-reload
sudo systemctl start demo-http
sudo systemctl status demo-http
curl http://localhost:9000
\end{lstlisting}

    \item Aktifkan auto-start lalu bersihkan setelah selesai:
\begin{lstlisting}[language=bash, caption={Enable dan pembersihan layanan kustom}]
sudo systemctl enable demo-http
sudo systemctl disable --now demo-http
sudo rm /etc/systemd/system/demo-http.service
sudo systemctl daemon-reload
\end{lstlisting}
\end{enumerate}

%% ========================================================================
\section{Rangkuman}
\label{sec:service-summary}
%% ========================================================================

Bab ini membahas pokok-pokok penting berikut:

\begin{description}
    \item[Sistem Init] Linux modern menggunakan systemd sebagai init system utama,
        pengganti pendekatan lama seperti SysV init dan Upstart.
    \item[systemctl] Perintah inti untuk start, stop, restart, reload, enable,
        disable, dan inspeksi layanan.
    \item[Unit File] Konfigurasi layanan dideklarasikan di file unit dengan blok
        \texttt{[Unit]}, \texttt{[Service]}, dan \texttt{[Install]}.
    \item[Dependensi] Relasi seperti \texttt{After=}, \texttt{Wants=}, dan
        \texttt{Requires=} mengatur urutan dan keterkaitan layanan.
    \item[Logging] \cmd{journalctl} dan \cmd{systemctl status} adalah alat utama
        untuk observasi dan diagnosis layanan.
    \item[Layanan Jaringan] SSH, Apache, dan MySQL adalah contoh penting layanan
        yang perlu dikelola, divalidasi, dan dimonitor.
    \item[Automatisasi] systemd mendukung startup otomatis, hook layanan, dan
        timer untuk maintenance terjadwal.
    \item[Troubleshooting] Diagnosis layanan biasanya dimulai dari status, log,
        port, path binary, permission, dan dependensi eksternal.
\end{description}

%% ========================================================================
\section{Latihan}
\label{sec:service-exercises}
%% ========================================================================

\begin{exercisebox}[12.1]
\begin{enumerate}
    \item Apa perbedaan praktis antara \cmd{restart} dan \cmd{reload} pada layanan?
        Berikan contoh layanan yang lebih aman di-\cmd{reload} daripada di-restart.
    \item Mengapa \cmd{mask} lebih kuat daripada \cmd{disable}? Dalam skenario apa
        Anda memilih \cmd{mask}?
    \item Jelaskan perbedaan fungsi \texttt{After=} dan \texttt{Requires=} dalam
        unit file systemd.
\end{enumerate}
\end{exercisebox}

\begin{exercisebox}[12.2]
\begin{enumerate}
    \item Buat rancangan unit file untuk aplikasi Python sederhana yang harus:
        berjalan sebagai user \texttt{www-data}, start setelah jaringan aktif,
        dan otomatis restart saat gagal.
    \item Sebuah layanan gagal start dengan pesan bahwa port 3306 sudah digunakan.
        Tuliskan langkah diagnostik lengkap yang akan Anda lakukan.
    \item Bandingkan kapan Anda lebih tepat memakai cron, dan kapan lebih tepat
        memakai systemd timer.
\end{enumerate}
\end{exercisebox}

%% ========================================================================
\section{Referensi Tambahan}
\label{sec:service-references}
%% ========================================================================

\begin{itemize}
    \item \textbf{Halaman manual}:
        \cmd{man systemctl}, \cmd{man systemd}, \cmd{man systemd.service},
        \cmd{man journalctl}, \cmd{man systemd.timer}
    \item \textbf{Freedesktop systemd documentation}:
        \href{https://www.freedesktop.org/software/systemd/man/latest/}{freedesktop.org/software/systemd/man/latest/}
    \item \textbf{Ubuntu Server Guide: Service Management}:
        \href{https://documentation.ubuntu.com/server/}{documentation.ubuntu.com/server/}
    \item \textbf{DigitalOcean tutorial systemd}:
        \href{https://www.digitalocean.com/community/tutorials/systemd-essentials-working-with-services-units-and-the-journal}{digitalocean.com/community/tutorials/systemd-essentials-working-with-services-units-and-the-journal}
\end{itemize}
